\let\negmedspace\undefined
\let\negthickspace\undefined
\documentclass[journal]{IEEEtran}
\usepackage[a5paper, margin=10mm, onecolumn]{geometry}
\usepackage{lmodern} 
\usepackage{tfrupee} 
\setlength{\headheight}{1cm}
\setlength{\headsep}{0mm}   

\usepackage{gvv-book}
\usepackage{gvv}
\usepackage{cite}
\usepackage{amsmath,amssymb,amsfonts,amsthm}
\usepackage{algorithmic}
\usepackage{graphicx}
\usepackage{textcomp}
\usepackage{xcolor}
\usepackage{txfonts}
\usepackage{listings}
\usepackage{enumitem}
\usepackage{mathtools}
\usepackage{gensymb}
\usepackage{comment}
\usepackage[breaklinks=true]{hyperref}
\usepackage{tkz-euclide} 
\usepackage{listings}  
\def\inputGnumericTable{}                                 
\usepackage[latin1]{inputenc}                                
\usepackage{color}                                            
\usepackage{array}                                            
\usepackage{longtable}                                       
\usepackage{calc}                                             
\usepackage{multirow}                                         
\usepackage{hhline}                                           
\usepackage{ifthen}                                           
\usepackage{lscape}
\usepackage{xparse}

\bibliographystyle{IEEEtran}

\title{1.2.40}
\author{EE25BTECH11062 - Vivek K Kumar}

\begin{document}
\maketitle

\renewcommand{\thefigure}{\theenumi}
\renewcommand{\thetable}{\theenumi}

\numberwithin{equation}{enumi}
\numberwithin{figure}{enumi} 

\textbf{Question}:\\
\begin{align}
    \lim\limits_{x\to0} \frac{e^x-\brak{1+x+\frac{x^2}{2}}}{x^3} = 
\end{align}
\begin{enumerate}
    \item $0$
    \item $1/6$
    \item $1/3$
    \item $1$
\end{enumerate}
\textbf{Solution: }

\textbf{Method 1 - Taylor series}\\
Taylor series expansion of $e^x$
\begin{align}
    e^x = \sum\limits_{n=0}^\infty \frac{x^n}{n!} = 1+x+\frac{x^2}{2}+\frac{x^3}{6}+\frac{x^4}{24} \dots
\end{align}

Then the limit becomes
\begin{align}
    &\lim\limits_{x\to0}\frac{\brak{1+x+\frac{x^2}{2}+\frac{x^3}{6}+\frac{x^4}{24} \dots} - \brak{1+x+\frac{x^2}{2}}}{x^3} \\
    &= \lim\limits_{x\to0} \frac{\frac{x^3}{6} + \frac{x^4}{24} + \dots}{x^3} \\
    &= \lim\limits_{x\to0} \frac{1}{6} + \frac{x}{24} + \frac{x^2}{120}\dots \\
   & = \boxed{\frac{1}{6}}
\end{align}
\begin{center}
$\boxed{\mathbf{ \lim\limits_{x\to0} \frac{e^x-\brak{1+x+\frac{x^2}{2}}}{x^3} = \frac{1}{6}}}$
\end{center}

\textbf{Method 2 - L'Hospital Rule}
Since numerator and denominator tends to 0, we can apply L'Hospital rule. 

\begin{align}
    &\lim\limits_{x\to0} \frac{e^x-\brak{1+x+\frac{x^2}{2}}}{x^3} \\
    &= \lim\limits_{x\to0}\frac{\frac{d}{dx}\brak{e^x - \brak{1+x+\frac{x^2}{2}}}}{\frac{d}{dx}(x^3)} \\
    &= \lim\limits_{x\to0} \frac{e^x - \brak{1+x}}{3x^2}    
\end{align}

Applying L'hospital again 

\begin{align}
    & \lim\limits_{x\to0} \frac{e^x - \brak{1+x}}{3x^2} \\
    &= \lim\limits_{x\to0}\frac{\frac{d}{dx}\brak{e^x - \brak{1+x}}}{\frac{d}{dx}(3x^2)} \\
    &= \lim\limits_{x\to0} \frac{e^x - 1}{6x} \\
\end{align}

Applying L'hospital one last time
\begin{align}
    & \lim\limits_{x\to0} \frac{e^x - 1}{6x} \\
    &= \lim\limits_{x\to0}\frac{\frac{d}{dx}\brak{e^x - 1}}{\frac{d}{dx}(6x)} \\
    &= \lim\limits_{x\to0} \frac{e^x}{6} \\
    & = \frac{e^0}{6} = \boxed{\frac{1}{6}}
\end{align}

\begin{center}
$\boxed{\textbf{Option 2 is the correct answer}}$
\end{center}


\end{document}  

